%! Introduction to Eigenvalues — Unit 6, Lesson 6.1 — Homework (Answer Key)
\documentclass[10pt]{article}
\usepackage{linalg-article}
\usepackage{linalg-key}   % pulls in -boxes; do NOT also load -boxes
\newcommand{\TallMath}[1]{$\displaystyle #1\rule[-1.4em]{0pt}{3.2em}$}
\newcommand{\xx}{\mathbf{x}}
\newcommand{\zero}{\mathbf{0}}

\begin{document}

\pageheader{Unit 6, Lesson 6.1}{Homework --- Answer Key}
\namedateperiod

\vspace{0.12in}

\begin{notesbox}{Practice}
The method is two steps. \textbf{Step 1:} solve $\det(A-\lambda I)=0$ for the \textbf{eigenvalues}.
\textbf{Step 2:} for each $\lambda$, solve $(A-\lambda I)\xx=\zero$ for an \textbf{eigenvector}. Check
with the trace ($=\lambda_1+\lambda_2$) and determinant ($=\lambda_1\lambda_2$).

\begin{enumerate}[leftmargin=1.9em, itemsep=12pt]
  \item \textbf{Full method.} For $A=\begin{bmatrix} 5 & 2 \\ 2 & 5 \end{bmatrix}$, find both
        eigenvalues, then an eigenvector for each.
        \ansline{$\det(A-\lambda I)=(5-\lambda)^2-4=\lambda^2-10\lambda+21=(\lambda-3)(\lambda-7)\Rightarrow\lambda=3,7$.}
        \ansline{$\lambda=7$: $\begin{bsmallmatrix}-2&2\\2&-2\end{bsmallmatrix}\Rightarrow\begin{bsmallmatrix}1\\1\end{bsmallmatrix}$; \quad $\lambda=3$: $\begin{bsmallmatrix}2&2\\2&2\end{bsmallmatrix}\Rightarrow\begin{bsmallmatrix}1\\-1\end{bsmallmatrix}$.}
  \item \textbf{Read the eigenvalues off.} State the eigenvalues of each matrix (no full work needed ---
        look at the diagonal), and say why.
    \begin{enumerate}[label=(\alph*), leftmargin=1.8em, itemsep=4pt]
      \item $A=\begin{bmatrix} 6 & 0 \\ 0 & 2 \end{bmatrix}$ \hfill \ans{$\lambda=6,2$ (diagonal).}
      \item $A=\begin{bmatrix} 4 & 5 \\ 0 & 1 \end{bmatrix}$ \;(triangular) \hfill \ans{$\lambda=4,1$ (diagonal).}
    \end{enumerate}
    \ansline{For a diagonal or triangular matrix $\det(A-\lambda I)=(4-\lambda)(1-\lambda)$ is already factored, so the eigenvalues are the diagonal entries.}
  \item \textbf{A negative eigenvalue.} For $A=\begin{bmatrix} 1 & 2 \\ 2 & 1 \end{bmatrix}$, find both
        eigenvalues and an eigenvector for each. What does the \emph{negative} eigenvalue do to its
        direction?
        \ansline{$\det(A-\lambda I)=(1-\lambda)^2-4=\lambda^2-2\lambda-3=(\lambda-3)(\lambda+1)\Rightarrow\lambda=3,-1$.}
        \ansline{$\lambda=3$: $\begin{bsmallmatrix}-2&2\\2&-2\end{bsmallmatrix}\Rightarrow\begin{bsmallmatrix}1\\1\end{bsmallmatrix}$; \quad $\lambda=-1$: $\begin{bsmallmatrix}2&2\\2&2\end{bsmallmatrix}\Rightarrow\begin{bsmallmatrix}1\\-1\end{bsmallmatrix}$. The negative eigenvalue \emph{reverses} (flips) its direction while keeping the same line.}
  \item \textbf{Trace \& determinant check.} For $A=\begin{bmatrix} 4 & 1 \\ 2 & 3 \end{bmatrix}$ (the
        warm-up matrix), find both eigenvalues. Then verify that their \emph{sum} equals the trace and
        their \emph{product} equals $\det A$.
        \ansline{$\det(A-\lambda I)=(4-\lambda)(3-\lambda)-2=\lambda^2-7\lambda+10=(\lambda-2)(\lambda-5)\Rightarrow\lambda=2,5$. Sum $2+5=7=$ trace $(4+3)$; product $2\cdot5=10=\det A\ (12-2)$. Both check.}
  \item \textbf{Explain.} Why does solving $\det(A-\lambda I)=0$ give the eigenvalues, and why must an
        eigenvector be \emph{nonzero}?
        \ansline{$\lambda$ is an eigenvalue exactly when $(A-\lambda I)\xx=\zero$ has a \emph{nonzero} solution, which needs $A-\lambda I$ singular --- i.e.\ $\det(A-\lambda I)=0$. The zero vector solves $(A-\lambda I)\xx=\zero$ for \emph{every} $\lambda$, so allowing it would give no information; that is why eigenvectors must be nonzero.}
\end{enumerate}
\end{notesbox}

\begin{extensionbox}
\textbf{Eigenvalues of a $3\times3$ triangular matrix.} The method scales up: for a triangular matrix,
$\det(A-\lambda I)$ is the product down the diagonal. Find the eigenvalues of
\[
  A=\begin{bmatrix} 2 & 1 & 7 \\ 0 & 3 & 4 \\ 0 & 0 & 5 \end{bmatrix}.
\]
(\emph{Hint:} $\det(A-\lambda I)=(2-\lambda)(3-\lambda)(5-\lambda)$.)
\ansline{$(2-\lambda)(3-\lambda)(5-\lambda)=0\Rightarrow\lambda=2,3,5$ --- again the diagonal entries. (For triangular matrices of any size, the eigenvalues are the diagonal.)}
\end{extensionbox}

\begin{spiralbox}
\textbf{Coming next --- Lesson 6.2: Diagonalizing a Matrix.} You can now \emph{find} a matrix's
eigenvalues and eigenvectors. Next we line up the eigenvectors into a matrix and use them to
\emph{rebuild} $A$ as $A=X\Lambda X^{-1}$ --- turning powers, and hard matrix problems, into easy work
on the numbers $\lambda$.
\end{spiralbox}

\begin{teachernote}
Item~1 is the full two-step method on a clean symmetric matrix ($\lambda=3,7$). Item~2 is the
diagonal/triangular shortcut (eigenvalues sit on the diagonal --- a Lesson~6.2 foreshadow). Item~3
introduces a \emph{negative} eigenvalue (a flip); item~4 practices the trace/determinant check and
calls back to the warm-up matrix; item~5 is the target justification (singular $\Rightarrow\det=0$;
why $\xx\ne\zero$). The extension shows the method generalizing to $3\times3$ triangular. Most common
slips: sign errors forming $A-\lambda I$, dropping the ``$-bc$'' term, and reporting eigenvalues
without solving for an eigenvector.
\end{teachernote}

\end{document}
